% GENERATED FILE. Edit canonical lessons, then run npm run generate:books.
% canonical-source-hash: fnv1a64:9329d21ff1404c08
% canonical-lessons: IT-C307-nascita, IT-C307-gemelli, IT-C307-primogenito, IT-C307-svezzare, IT-C307-ciuccio

\chapter{Benvenuto al mondo}
\label{ch:it-benvenuto-al-mondo}
\hlchaptermodality{307}

\begin{chapteropening}
\textbf{By the end of this chapter:} \emph{I can talk about a birth: twins, the first-born, weaning and the dummy.}

Complete the last lesson of chapter 307: I can talk about a birth: twins, the first-born, weaning and the dummy.
\end{chapteropening}

\section[la nascita]{la nascita — the birth}

\label{lesson:IT-C307-nascita}



\begin{quote}
Before the new one: say the Italian for missing, then the Italian for to survive.
\end{quote}

\subsection*{You'll want to know: la nascita}
\textbf{la nascita} — \textquotedblleft{}the birth\textquotedblright{}.

\begin{quote}
La nascita di mia figlia è stata il giorno più bello della mia vita.
\end{quote}

\begin{quote}
\emph{My daughter's birth was the best day of my life.}
\end{quote}

\emph{La data di nascita} — the date of birth.

\begin{cousinweb}
\emph{nascita} is built on \emph{nascere} \textquotedblleft{}to be born\textquotedblright{}, which is inherited from Latin \emph{nāscī} (Lewis \& Short).
\end{cousinweb}

\begin{grammarlens}[title={What you've built}]
One more word for telling, explaining and describing at B1.
\end{grammarlens}

\subsection*{Your turn}
\begin{itemize}
\raggedright
  \item \emph{Say it:} \emph{la nascita}
  \item \emph{Say it:} \emph{la nascita}, once more
  \item \emph{Say it:} a sentence of your own with \emph{la nascita}
  \item \emph{Recall:} say the Italian for missing, then the Italian for to survive, then say \emph{la nascita} again
\end{itemize}

\begin{tcolorbox}[breakable,colback=teal!4,colframe=teal!35!black,title={Before you move on}]
What does la nascita mean? (\textquotedblleft{}The birth\textquotedblright{}.) Say it once more.
\end{tcolorbox}
\section[i gemelli]{i gemelli — the twins}

\label{lesson:IT-C307-gemelli}



\begin{quote}
Before the new one: say the Italian for to survive, then the Italian for the birth.
\end{quote}

\subsection*{You'll want to know: i gemelli}
\textbf{i gemelli} — \textquotedblleft{}the twins\textquotedblright{}.

\begin{quote}
Mia sorella ha avuto due gemelli.
\end{quote}

\begin{quote}
\emph{My sister has had twins.}
\end{quote}

Singular \emph{il gemello}, \emph{la gemella}; \emph{le gemelle} for two girls.

\begin{cousinweb}
\emph{gemelli} is inherited from Latin \emph{gemellus} \textquotedblleft{}a twin\textquotedblright{}, a little \emph{geminus} \textquotedblleft{}twin-born\textquotedblright{} (Lewis \& Short).
\end{cousinweb}

\begin{grammarlens}[title={What you've built}]
One more word for telling, explaining and describing at B1.
\end{grammarlens}

\subsection*{Your turn}
\begin{itemize}
\raggedright
  \item \emph{Say it:} \emph{i gemelli}
  \item \emph{Say it:} \emph{i gemelli}, once more
  \item \emph{Say it:} a sentence of your own with \emph{i gemelli}
  \item \emph{Recall:} say the Italian for to survive, then the Italian for the birth, then say \emph{i gemelli} again
\end{itemize}

\begin{tcolorbox}[breakable,colback=teal!4,colframe=teal!35!black,title={Before you move on}]
What does i gemelli mean? (\textquotedblleft{}The twins\textquotedblright{}.) Say it once more.
\end{tcolorbox}
\section[il primogenito]{il primogenito — the first-born child}

\label{lesson:IT-C307-primogenito}



\begin{quote}
Before the new one: say the Italian for the birth, then the Italian for the twins.
\end{quote}

\subsection*{You'll want to know: il primogenito}
\textbf{il primogenito} — \textquotedblleft{}the first-born child\textquotedblright{}.

\begin{quote}
Il primogenito si chiama come il nonno.
\end{quote}

\begin{quote}
\emph{The first-born is named after his grandfather.}
\end{quote}

A daughter is \emph{la primogenita}.

\begin{cousinweb}
\emph{primogenito} is a learned borrowing from Latin \emph{prīmōgenitus} \textquotedblleft{}first-born\textquotedblright{} (EtymDB), from \emph{prīmus} \textquotedblleft{}first\textquotedblright{} and \emph{genitus} \textquotedblleft{}born\textquotedblright{}.
\end{cousinweb}

\begin{grammarlens}[title={What you've built}]
One more word for telling, explaining and describing at B1.
\end{grammarlens}

\subsection*{Your turn}
\begin{itemize}
\raggedright
  \item \emph{Say it:} \emph{il primogenito}
  \item \emph{Say it:} \emph{il primogenito}, once more
  \item \emph{Say it:} a sentence of your own with \emph{il primogenito}
  \item \emph{Recall:} say the Italian for the birth, then the Italian for the twins, then say \emph{il primogenito} again
\end{itemize}

\begin{tcolorbox}[breakable,colback=teal!4,colframe=teal!35!black,title={Before you move on}]
What does il primogenito mean? (\textquotedblleft{}The first-born child\textquotedblright{}.) Say it once more.
\end{tcolorbox}
\section[svezzare]{svezzare — to wean}

\label{lesson:IT-C307-svezzare}



\begin{quote}
Before the new one: say the Italian for the twins, then the Italian for the first-born child.
\end{quote}

\subsection*{You'll want to know: svezzare}
\textbf{svezzare} — \textquotedblleft{}to wean\textquotedblright{}.

\begin{quote}
Abbiamo svezzato il bambino a sei mesi.
\end{quote}

\begin{quote}
\emph{We weaned the baby at six months.}
\end{quote}

Literally to take out of a habit: \emph{vezzo} is an old word for habit.

\begin{cousinweb}
\emph{svezzare} is built on \emph{vezzo} \textquotedblleft{}habit\textquotedblright{}, from Latin \emph{vitium} \textquotedblleft{}a fault\textquotedblright{} (Lewis \& Short), with \emph{s-} \textquotedblleft{}out of\textquotedblright{}.
\end{cousinweb}

\begin{grammarlens}[title={What you've built}]
One more word for telling, explaining and describing at B1.
\end{grammarlens}

\subsection*{Your turn}
\begin{itemize}
\raggedright
  \item \emph{Say it:} \emph{svezzare}
  \item \emph{Say it:} \emph{svezzare}, once more
  \item \emph{Say it:} a sentence of your own with \emph{svezzare}
  \item \emph{Recall:} say the Italian for the twins, then the Italian for the first-born child, then say \emph{svezzare} again
\end{itemize}

\begin{tcolorbox}[breakable,colback=teal!4,colframe=teal!35!black,title={Before you move on}]
What does svezzare mean? (\textquotedblleft{}To wean\textquotedblright{}.) Say it once more.
\end{tcolorbox}
\section[il ciuccio]{il ciuccio — the dummy, the pacifier}

\label{lesson:IT-C307-ciuccio}



\begin{quote}
Before the new one: say the Italian for the first-born child, then the Italian for to wean.
\end{quote}

\subsection*{You'll want to know: il ciuccio}
\textbf{il ciuccio} — \textquotedblleft{}the dummy, the pacifier\textquotedblright{}.

\begin{quote}
Il bambino non dorme senza ciuccio.
\end{quote}

\begin{quote}
\emph{The baby will not sleep without his dummy.}
\end{quote}

Informal and very common; in shops also \emph{il succhiotto}.

\begin{cousinweb}
\emph{ciuccio} is built on \emph{ciucciare} \textquotedblleft{}to suck\textquotedblright{}, which is usually explained as a word that imitates the sound of sucking (EtymDB has no entry).
\end{cousinweb}

\begin{grammarlens}[title={What you've built}]
One more word for telling, explaining and describing at B1.
\end{grammarlens}

\subsection*{Your turn}
\begin{itemize}
\raggedright
  \item \emph{Say it:} \emph{il ciuccio}
  \item \emph{Say it:} \emph{il ciuccio}, once more
  \item \emph{Say it:} a sentence of your own with \emph{il ciuccio}
  \item \emph{Recall:} say the Italian for the first-born child, then the Italian for to wean, then say \emph{il ciuccio} again
\end{itemize}

\begin{tcolorbox}[breakable,colback=teal!4,colframe=teal!35!black,title={Before you move on}]
What does il ciuccio mean? (\textquotedblleft{}The dummy, the pacifier\textquotedblright{}.) Say it once more.
\end{tcolorbox}
